% ============================================================================
% chapters/ch4_implementation.tex
% ============================================================================
\chapter{Implementation}
\label{ch:implementation}

% ============================================================================
\section{Repository Structure}
\label{sec:repo_structure}
% ============================================================================

The implementation is organised as an installable Python package rooted in
the \texttt{src/} directory.
The repository follows a layered architecture separating concerns into four
tiers: entities, components, core modules, and the pipeline orchestrator.

\begin{figure}[ht]
  \centering
\begin{lstlisting}[language={},basicstyle=\ttfamily\footnotesize,
                   numbers=none, backgroundcolor=\color{backcolour}]
MasterArbeit_MLops/
|-- run_pipeline.py          <- Single entry point (CLI)
|-- pyproject.toml           <- Package definition, tool config
|-- src/
|   |-- entity/
|   |   |-- config_entity.py   <- 14 typed config dataclasses
|   |   `-- artifact_entity.py <- 14 typed artifact dataclasses
|   |-- components/          <- One file per pipeline stage (14 files)
|   |-- pipeline/
|   |   `-- production_pipeline.py  <- Orchestrator (470 lines)
|   |-- domain_adaptation/
|   |   |-- tent.py          <- TENT entropy minimisation
|   |   `-- adabn.py         <- Adaptive Batch Normalisation
|   |-- api/
|   |   `-- app.py           <- FastAPI inference server
|   |-- mlflow_tracking.py   <- Experiment tracking wrapper
|   |-- trigger_policy.py    <- Trigger engine (812 lines)
|   `-- train.py             <- Training with CV (925 lines)
|-- tests/                   <- 18 test files, conftest.py
|-- config/                  <- YAML configuration files
|-- docker/                  <- Dockerfiles and docker-compose.yml
`-- .github/workflows/
    `-- ci-cd.yml            <- CI/CD workflow (6 jobs)
\end{lstlisting}
  \caption{Repository directory structure (abbreviated).}
  \label{fig:repo_structure}
\end{figure}

% ============================================================================
\section{Entity and Component Layers}
\label{sec:entity_component}
% ============================================================================

\subsection{Entity Layer}

The entity layer defines the contract between stages using Python dataclasses.
Each stage receives a typed configuration dataclass and returns a typed
artefact dataclass.

Configuration dataclasses (\texttt{config\_entity.py}, 344 lines) encode
all parameters a stage might need with sensible defaults.
A master \texttt{PipelineConfig} holds project-wide settings from which
stage-specific configs derive their paths.

Artefact dataclasses (\texttt{artifact\_entity.py}, 217 lines) are the typed
outputs of each stage.
A \texttt{PipelineResult} aggregator collects all artefacts from a run and
enables serialisation of the full execution record to JSON.

This typed contract eliminates the class of bug where a stage silently receives
an incorrect file path or missing configuration value.

\subsection{Component Layer}

Each of the 14 pipeline stages is implemented as a component class following
a uniform interface:

\begin{lstlisting}[caption={Component interface pattern.},
                   label={lst:component}]
class DataIngestionComponent:
    def __init__(self, config: DataIngestionConfig):
        self.config = config

    def run(self) -> DataIngestionArtifact:
        # delegate to core module
        result = sensor_data_pipeline.ingest(self.config)
        return DataIngestionArtifact(output_path=result.path)
\end{lstlisting}

Components are responsible for: reading configuration, invoking the core
module, wrapping the result in a typed artefact, and logging entry/exit
messages.
All algorithmic complexity resides in the core modules, which can be unit-tested
without instantiating the pipeline.

% ============================================================================
\section{Pipeline Runner and CLI}
\label{sec:cli}
% ============================================================================

\texttt{run\_pipeline.py} is the single entry point.
All pipeline configurations are expressed as command-line arguments so that
any run is fully reproducible from its recorded command:

\begin{lstlisting}[caption={CLI usage examples.}]
# Default: inference cycle (Stages 1-7)
python run_pipeline.py

# Trigger-reactive retraining with AdaBN+TENT
python run_pipeline.py --retrain --adapt adabn_tent

# Full pipeline with pseudo-labelling
python run_pipeline.py --retrain --adapt pseudo_label \
    --update-baseline --force-retrain

# Batch processing of 26 sessions
python batch_process_all_datasets.py
\end{lstlisting}

The runner selects stages based on flags and resolves missing upstream
artefacts via fallback to the most recent successful artefact directory.
This enables incremental execution: if Stage 3 artefacts already exist,
Stages 1--2 are skipped on a re-run.

% ============================================================================
\section{Tooling Stack}
\label{sec:tooling}
% ============================================================================

\subsection{DVC — Data Version Control}

DVC \citep{dvcteam2020} tracks every data file and model checkpoint through
pointer files committed to Git.
Each pointer stores the MD5 hash and storage path of the tracked binary.
\texttt{dvc pull} retrieves all tracked files at the version corresponding to
a given Git commit, ensuring that reproduction of any prior result requires
only \texttt{git checkout} followed by \texttt{dvc pull}.

Key tracked artefacts:
\begin{itemize}
  \item \texttt{data/raw/sensor\_fused\_50Hz.csv} — raw ingestion output
  \item \texttt{data/prepared/X\_train.npy}, \texttt{y\_train.npy} — training windows
  \item \texttt{models/pretrained/fine\_tuned\_model\_1dcnnbilstm.keras} — production model
  \item \texttt{models/normalized\_baseline.json} — reference distribution
\end{itemize}

\subsection{MLflow — Experiment Tracking and Model Registry}

MLflow \citep{mlflowteam2018} records all experiment metadata in a local
SQLite database under \texttt{mlruns/}.
The \texttt{MLflowTracker} class in \texttt{src/mlflow\_tracking.py} wraps
all MLflow API calls across the pipeline.

Four experiments are used:
\begin{description}
  \item[\texttt{har-training}] Training runs: hyperparameters,
        per-epoch metrics, model artefacts.
  \item[\texttt{har-retraining}] Adaptation runs: pre/post confidence,
        TENT rollback flag, entropy delta, canary results.
  \item[\texttt{har-monitoring}] Per-session monitoring: PSI scores,
        drift flags, z-scores, trigger decisions.
  \item[\texttt{har-evaluation}] Evaluation runs: F1 per class,
        ECE, calibration plot artefact.
\end{description}

Model promotion is controlled through the MLflow Model Registry.
A model version may be in one of three states: \texttt{Staging},
\texttt{Production}, or \texttt{Archived}.
A canary evaluation gate must be passed before promotion from Staging to
Production.

\subsection{Docker — Containerised Deployment}

Two Docker services are defined in \texttt{docker-compose.yml}:

\begin{description}
  \item[\texttt{har-mlflow}] MLflow tracking server on port 5000.
        Mounts \texttt{./mlruns} and \texttt{./models} for persistence.
  \item[\texttt{har-inference}] FastAPI inference server on port 8000.
        Built from \texttt{docker/Dockerfile.inference} using
        Python 3.11-slim base image.
        Sets \texttt{PYTHONPATH=/app:/app/src} and entrypoint
        \texttt{uvicorn src.api.app:app}.
\end{description}

\subsection{FastAPI — Inference API}

The inference server provides two primary endpoints:

\begin{table}[ht]
  \centering
  \caption{FastAPI endpoint specification.}
  \label{tab:api_endpoints}
  \begin{tabularx}{\linewidth}{llX}
    \toprule
    \textbf{Endpoint} & \textbf{Method} & \textbf{Description} \\
    \midrule
    \texttt{/api/predict}  & POST & Accept multipart CSV upload; return
      predictions, confidence, monitoring summary \\
    \texttt{/api/health}   & GET  & Health check; returns
      \texttt{200 OK} with model load status \\
    \texttt{/docs}         & GET  & Auto-generated Swagger UI \\
    \bottomrule
  \end{tabularx}
\end{table}

The handler applies the full preprocessing pipeline (unit detection $\to$
StandardScaler $\to$ windowing) before calling the model, and runs all three
monitoring layers on the resulting predictions before returning the response.

% ============================================================================
\section{CI/CD Workflow}
\label{sec:cicd}
% ============================================================================

The GitHub Actions workflow (\texttt{.github/workflows/ci-cd.yml}) runs on
every push to \texttt{main} and on a scheduled nightly trigger.
It is split into six jobs:

\begin{enumerate}
  \item \textbf{Lint}: \texttt{black}, \texttt{flake8}, \texttt{mypy}
  \item \textbf{Test-fast}: \texttt{pytest -m "not slow"} (pure Python, no TF)
  \item \textbf{Test-slow}: \texttt{pytest -m slow} (TF-dependent, longer timeout)
  \item \textbf{Build}: Docker image build for \texttt{har-inference}
  \item \textbf{Integration}: smoke test against the running container
          at \texttt{/api/health} and \texttt{/api/predict}
  \item \textbf{Push}: push to GHCR as \texttt{har-inference:latest}
\end{enumerate}

Fast and slow tests are separated to provide immediate feedback on small changes
without timing out on TensorFlow model loading (see \Cref{sec:engineering_challenges}).

% ============================================================================
\section{Engineering Challenges}
\label{sec:engineering_challenges}
% ============================================================================

This section summarises selected engineering problems encountered during
development.
A complete log with 21 entries, each backed by a Git commit reference, is
provided in \Cref{app:engineering_log}.

\subsection{Accelerometer Unit Mismatch}
\label{sec:unit_mismatch}

The most consequential problem was a sensor unit incompatibility:
production data exported from Garmin Connect was in milli-g, while the
pre-trained model was trained on data in m/s\textsuperscript{2}.
After applying the training StandardScaler to the production data, the
normalised feature values were approximately 100$\times$ out of range.
The model exhibited 14.5\,\% cross-user accuracy with high-confidence
wrong predictions --- below the random-chance baseline of 9.09\,\% for
\nclasses{} classes.

Root cause: the training dataset's $a_z$ mean was $-3.5$\,m/s\textsuperscript{2}
(approximately $-1g$ vertical); the production mean was $-1001.6$\,milli-g
(the same physical value, different unit).
The fix was the automatic unit-detection step described in \Cref{sec:dataset}.
Evidence: commits \texttt{c17f3dd}, \texttt{3a07e4f}, \texttt{1c403f8}.

\subsection{TENT Corrupts AdaBN Running Statistics}

When AdaBN and TENT are chained, TENT's training-mode forward passes update
BN running statistics, overwriting the target-domain statistics set by AdaBN.
The snapshot-and-restore mechanism described in \Cref{sec:tent_adabn}
resolved this.
Evidence: commit \texttt{f47a48b}.

\subsection{Pseudo-Label Source Data Scale Mismatch}

In the pseudo-label retraining loop, source-domain data was passed at raw scale
while pseudo-labelled target data was already scaled.
Mixing the two inflated in-loop validation accuracy but degraded production
performance.
Fix: scale source data with the production StandardScaler before mixing.
Evidence: commit \texttt{3fd3c00}.

\subsection{PSI Threshold Calibration}

Standard textbook PSI thresholds (0.10 / 0.25) were applied to the multi-channel
aggregated PSI.
These thresholds are correct for single-channel distributions but are far too
sensitive for a 24-channel aggregated score.
Empirically calibrated thresholds (0.75 / 1.50) were substituted.
Evidence: commit \texttt{bd8dc1e}.
